\section{Non-bonding MOs, disjointness and the sign of $K$ (Borden--Davidson)}\label{sec:huckel}

\begin{reviewclaim}
``Borden and Davidson proposed to classify the nonbonding molecular orbitals (NBMOs) $\varphi_a$ and
$\varphi_b$ as either disjoint or nondisjoint. If the two orbitals do not have common atoms and the
probability of two electrons occupying the same space is nil independent of their spins, they are
termed disjoint. Under such circumstances, the energy of an open-shell singlet can become comparable to
or even lower than that of a triplet.'' (p.~443) Examples: CBD, TME (disjoint, singlet); MX, TriMB
(nondisjoint, high spin) (p.~443--444, Figure~3).
\end{reviewclaim}

\subsection{Hückel NBMOs: the zero-sum rule}
Hückel Hamiltonian for a conjugated $\pi$ system with adjacency matrix $\mathbf A$ ($A_{ij}=1$ if atoms
$i,j$ are bonded): $\mathbf H=\alpha\mathbf 1+\beta\mathbf A$. Writing $E=\alpha+x\beta$, the secular equation
is $\mathbf A\mathbf c=x\mathbf c$.
\begin{enumerate}
  \item An NBMO has $E=\alpha$, i.e.\ $x=0$, so $\mathbf A\mathbf c=\mathbf 0$.
  \item Row $i$ of this equation reads $\sum_{j\in\mathcal N(i)}c_j=0$. \textbf{Zero-sum rule:} for every atom, the
        NBMO coefficients of its neighbours sum to zero.
\end{enumerate}

\subsection{Alternant hydrocarbons: pairing theorem and NBMO count}
If the atoms split into ``starred'' ($\mathcal A$, $n_A$ atoms) and ``unstarred'' ($\mathcal B$, $n_B$ atoms) sets with
bonds only between the sets (a bipartite graph), then in that ordering
$\mathbf A=\begin{pmatrix}\mathbf0&\mathbf B\\\mathbf B^{\mathsf T}&\mathbf0\end{pmatrix}$, with $\mathbf B$ of size $n_A\times n_B$.
\begin{enumerate}
  \item \emph{Pairing theorem.} If $\mathbf A\binom{\mathbf u}{\mathbf v}=x\binom{\mathbf u}{\mathbf v}$, i.e.\
        $\mathbf B\mathbf v=x\mathbf u$ and $\mathbf B^{\mathsf T}\mathbf u=x\mathbf v$, then
        $\mathbf A\binom{\mathbf u}{-\mathbf v}=\binom{-\mathbf B\mathbf v}{\mathbf B^{\mathsf T}\mathbf u}=\binom{-x\mathbf u}{x\mathbf v}=-x\binom{\mathbf u}{-\mathbf v}$.
        Levels come in $\pm x$ pairs, with coefficients equal in magnitude and the sign flipped on one set.
  \item \emph{Lower bound from the sublattice imbalance.} $\operatorname{rank}\mathbf A=2\operatorname{rank}\mathbf B\le2\min(n_A,n_B)$.
        Hence the number of NBMOs (the nullity) is $N-\operatorname{rank}\mathbf A\ge N-2\min(n_A,n_B)=|n_A-n_B|$.
  \item \emph{Longuet-Higgins bound.} $\operatorname{rank}\mathbf B$ cannot exceed the largest number of
        non-attacking nonzero entries of $\mathbf B$. That is the size $T$ of a maximum matching, which equals
        the number of double bonds in the best Kekulé structure. So $\#\text{NBMO}\ge N-2T\ \ (\ge|n_A-n_B|)$.
        \textbf{This is a lower bound, not an equality}: square cyclobutadiene has $T=2$, so $N-2T=0$,
        yet it has two NBMOs (its $\pm x$ pairs include an accidental $x=0$ pair).
\end{enumerate}

\subsection{Which NBMO pairs are disjoint? Explicit constructions}
\paragraph{TME (atoms 1,2 central; 3,4 on C1; 5,6 on C2).} Zero-sum at atoms 3 and 4 (single neighbour 1)
gives $c_1=0$. At atoms 5 and 6 it gives $c_2=0$. At atom 1: $c_2+c_3+c_4=0$, so $c_4=-c_3$. At atom 2:
$c_5=-c_6$. The NBMO space is spanned by
\begin{equation}
  \varphi_a=\tfrac1{\sqrt2}(\chi_3-\chi_4),\qquad\varphi_b=\tfrac1{\sqrt2}(\chi_5-\chi_6),
\end{equation}
which share no atom: \textbf{disjoint} (two allyl-type NBMOs on separate halves).

\paragraph{Square CBD (ring 1-2-3-4).} Zero-sum at atom 1: $c_2+c_4=0$; at atom 2: $c_1+c_3=0$. So
$\varphi_a=\tfrac1{\sqrt2}(\chi_1-\chi_3)$ and $\varphi_b=\tfrac1{\sqrt2}(\chi_2-\chi_4)$: \textbf{disjoint}.

\paragraph{TMM (centre 0, terminals 1,2,3).} $c_0=0$ (zero-sum at a terminal), and $c_1+c_2+c_3=0$ (zero-sum
at the centre). Every nonzero vector in this plane has at least two nonzero entries (one nonzero entry
cannot sum to zero). Two such vectors with disjoint supports would need at least $2+2=4$ terminal
atoms, but there are only 3. So no basis of the NBMO space is disjoint: \textbf{nondisjoint}.

\subsection{Why disjointness kills $K_{ab}$: zero differential overlap}
Expand $\varphi_a=\sum_\mu a_\mu\chi_\mu$, $\varphi_b=\sum_\nu b_\nu\chi_\nu$ in atomic $p_\pi$ orbitals:
\begin{enumerate}
  \item $K_{ab}=\eri abab=\sum_{\mu\nu\lambda\sigma}a_\mu b_\nu a_\lambda b_\sigma\eri{\mu}{\nu}{\lambda}{\sigma}$.
  \item ZDO approximation: $\eri\mu\nu\lambda\sigma=\delta_{\mu\nu}\delta_{\lambda\sigma}\gamma_{\mu\lambda}$, with $\gamma_{\mu\lambda}$ the
        Coulomb repulsion between charges on atoms $\mu$ and $\lambda$.
  \item Then $K_{ab}=\sum_{\mu\lambda}(a_\mu b_\mu)(a_\lambda b_\lambda)\gamma_{\mu\lambda}$.
  \item \emph{Disjoint} $\Rightarrow$ $a_\mu b_\mu=0$ for every atom $\Rightarrow$ $K_{ab}=0$ exactly (within ZDO).
        By \cref{eq:EOSS}, $E_{\text{OSS}}-E_T=2K_{ab}=0$: the first-order Hund preference vanishes, and
        higher-order effects (spin polarisation, \cref{sec:spinpol}) decide the ground state.
  \item \emph{Nondisjoint} $\Rightarrow$ some $a_\mu b_\mu\neq0$. Since $\gamma$ is a positive-definite Coulomb
        matrix (\cref{sec:holes}), $K_{ab}>0$ and the triplet is favoured at first order.
  \item In the Hubbard limit ($\gamma_{\mu\lambda}=U\delta_{\mu\lambda}$): $K_{ab}=U\sum_\mu a_\mu^2b_\mu^2$.
\end{enumerate}

\begin{verifybox}
T8 (\texttt{t06\_t09\_models.py}): Hückel spectra of CBD, TMM, TME, MX, TriMB, 1,2,4,5-TetraMB and
$p$-xylylene. Pairing theorem, zero-sum rule and $\#\text{NBMO}\ge N-2T\ge|n_A-n_B|$ hold in every case;
CBD shows the strict inequality. Minimising $\sum_\mu a_\mu^2b_\mu^2$ over all rotations of the degenerate
NBMO pair gives $\approx0$ for CBD, TME and TetraMB (disjoint), and $1/6$ (TMM) and $1/16$ (MX) (nondisjoint).
T9b (\texttt{t09b\_t15\_correlation.py}): in the Hubbard model at small $U$ the triplet--singlet gap of TMM
and MX is linear in $U$ (fitted exponent 0.97, 0.96). For TMM $-0.01643$ at $U=0.05t$ matches
$-2U\cdot\tfrac16=-0.0167$, and MX matches $-2U\cdot\tfrac1{16}$. For disjoint CBD and TME the gap is
$\propto U^{1.97},U^{1.95}$, i.e.\ zero at first order, as derived. \textbf{PASS}.
\end{verifybox}
\begin{sourcebox}
Borden \& Davidson, \emph{J.\ Am.\ Chem.\ Soc.}\ \textbf{99}, 4587 (1977), doi:10.1021/ja00456a010: not accessed;
definition as quoted in the review. Longuet-Higgins, \emph{J.\ Chem.\ Phys.}\ \textbf{18}, 265 (1950): the matching
bound as cited in arXiv:1906.08544. The proofs above are self-contained.
\end{sourcebox}
